\chapter{Études de cas de crise}\label{ch:m1:stress-case-studies}

À 14:45:28 le 6 mai 2010, le \omterm{def:m1:index-and-single-stock-futures-worldwide:ssf}{contrat à terme sur actions} le plus liquide du monde a cessé
de négocier pendant cinq secondes. Dans les quatre minutes et demie précédentes, il
avait chuté de plus de 5~\%~; les ordres d'achat en attente, supérieurs à cent mille
contrats ce matin-là, étaient tombés à environ mille. Quand la négociation a repris, le
cours s'est retourné et, en moins de vingt minutes, l'essentiel de la baisse avait été
effacé, sauf pour trois cents actions et fonds où, entre-temps, vingt mille transactions
s'étaient exécutées à des cours éloignés de plus de 60~\% de leur niveau antérieur~: un
cent, ou cent mille dollars. Rien dans l'économie n'a changé entre 14:32 et 15:00.
Chaque chapitre de ce livre décrit un mécanisme qui fonctionne un jour ordinaire.
Celui-ci examine quatre journées où les mécanismes ont atteint leurs limites, à partir
des rapports officiels, et pose à chaque fois les mêmes trois questions~: quel flux
est arrivé, qui était censé l'absorber, et pourquoi s'est-il retiré~?

\section{6 mai 2010}

\begin{definition}[Flash crash et cotation fictive]\label{def:m1:stress-case-studies:flash}
Un \emph{flash crash}\index{flash crash} est une chute suivie d'un rebond des cours, d'une
ampleur normalement observée sur plusieurs jours, en quelques minutes et sans nouvelle
correspondante. Une \emph{cotation fictive}\index{cotation fictive} est une cotation à un cours très éloigné du marché
(un cours acheteur d'un cent, un cours vendeur de 100\,000~\$), saisie par un \omterm{def:m1:what-a-trading-firm-does:market-maker}{teneur de marché} ou pour son compte afin de
satisfaire à une \omterm{def:m1:options-market-structure:mm}{obligation de cotation} continue sans avoir l'intention de négocier.
\end{definition}

Le rapport conjoint des deux régulateurs américains retrace la séquence. Le marché
était déjà en baisse et nerveux. À 14:32, un ensemble de fonds communs de placement a commencé à vendre
75\,000 contrats E-mini, soit environ 4.1~milliards de dollars, en tant que couverture, au moyen
d'un algorithme configuré pour vendre 9~\% du volume de la minute précédente, «~sans
considération de prix ni de temps~». Un programme comparable, mené par le même trader
avec des algorithmes qui tenaient compte du prix et du temps, avait auparavant mis plus de
cinq heures pour écouler ses premiers 75\,000 contrats~; celui-ci s'est achevé en vingt
minutes.

Les sociétés de haute fréquence et les autres intermédiaires ont acheté les premiers
contrats, puis, une fois atteintes leurs limites d'inventaire, les ont revendus. Entre
14:45:13 et 14:45:27, ils ont échangé plus de 27\,000 contrats, soit 49~\% du volume, tout en achetant
environ 200 contrats nets~: c'est l'effet «~patate chaude~» du rapport. Le volume ainsi engendré
a incité l'algorithme à vendre plus vite. À 14:45:28, la profondeur côté acheteur était tombée sous 1\,050
contrats, soit moins de 1~\% de son niveau du début de journée. La fonction de
logique d'arrêt de la bourse a alors suspendu la négociation pendant cinq secondes, et pendant cette
pause, les acheteurs sont revenus.

Sur le marché des actions, ce même retrait a eu une issue différente. Les \omterm{def:m1:what-a-trading-firm-does:market-maker}{teneurs de
marché} qui ont retiré leurs cotations n'ont laissé derrière eux que leurs \omterm{def:m1:stress-case-studies:flash}{cotations
fictives}~; les ordres au marché et les ordres stop, convertis en ordres au marché, se sont
exécutés contre elles. Les bourses ont ensuite annulé les transactions situées à plus
de 60~\% des cours de 14:40, un seuil fixé après coup, qui fut la dernière des
incertitudes de la journée pour quiconque avait acheté pendant la baisse.

\begin{omfigure}
\begin{tikzpicture}
\begin{axis}[width=11cm,height=5.2cm,
  xlabel={minutes depuis le début du programme},ylabel={variation du cours (\%)},
  tick label style={font=\small},label style={font=\small},
  xmin=1,xmax=40,ymin=-11,ymax=0.5,grid=major,grid style={gray!25},
  legend columns=3,legend style={font=\footnotesize,at={(0.5,-0.3)},anchor=north,draw=none,fill=none,/tikz/every even column/.append style={column sep=8pt}},
  table/col sep=comma]
\addplot[omExo,very thick,dotted] table[x=minute,y=none]{figdata/markets-1/31-stress-case-studies/paths.csv};
\addplot[omThm,very thick] table[x=minute,y=feedback]{figdata/markets-1/31-stress-case-studies/paths.csv};
\addplot[omDef,very thick,dashed] table[x=minute,y=pause]{figdata/markets-1/31-stress-case-studies/paths.csv};
\legend{un carnet fixe,la profondeur se retire et le volume tourne,la même chose avec une pause à $-5\%$}
\end{axis}
\end{tikzpicture}

\omcaption{Un modèle jouet du mécanisme, calibré sur rien~: un programme
vend 9~\% du volume de la dernière minute~; la profondeur se réduit, et le
tourbillon des intermédiaires fait monter le volume lorsque le cours est sous sa moyenne récente.
Avec un carnet fixe, les mêmes 75\,000 contrats coûtent 2.2~\%~; avec les deux
rétroactions, 9.8~\%, dont 5.9 points dans les quatre dernières minutes~; une pause qui
permet aux acheteurs de se réancrer arrête la baisse à 7.0~\%. Le modèle ne comporte pas de rebond. Données~: la simulation du tutoriel.}\label{fig:m1:stress-case-studies:paths}
\end{omfigure}

\section{24 août 2015}

\begin{definition}[Limit Up--Limit Down (LULD) et coupe-circuit de marché]\label{def:m1:stress-case-studies:luld}
\emph{Limit Up--Limit Down (LULD)}\index{Limit Up--Limit Down (LULD)} est le mécanisme américain,
introduit après 2010, qui empêche la négociation d'un titre en dehors d'une bande
autour d'un cours de référence mobile et suspend la négociation si le marché reste
au niveau de la bande. Un \emph{coupe-circuit de marché}\index{coupe-circuit de marché}
suspend toute la négociation des actions lorsqu'un indice large baisse, dans la
journée, de pourcentages fixés (voir au \cref{ch:m1:exchange-traded-funds}).
\end{definition}

\begin{dated}{2026-09}{Les règles de bande, telles que les décrivent les services du régulateur}\label{dat:m1:stress-case-studies:luld}
Le cours de référence est la moyenne des cours des transactions des cinq minutes
précédentes~; un nouveau cours de référence ne prend effet que s'il diffère d'au moins 1~\% du
précédent et s'il se maintient pendant au moins 30 secondes. Les bandes sont de 5~\%, 10~\% ou 20~\% de part et d'autre,
selon la catégorie du titre et son cours, doublées de 09:30 à 09:45 et de 15:35 à
16:00, et multipliées par le ratio de levier pour les produits à effet de levier. Lorsque le
meilleur cours vendeur atteint la bande inférieure, ou le meilleur cours acheteur la bande supérieure, le titre est
en état de limite~; si cet état dure 15 secondes, la bourse de référence déclare une
pause d'au moins cinq minutes, sur toutes les plateformes.
\end{dated}

Le lundi 24 août 2015, après des baisses en Asie, le E-mini a atteint sa
limite nocturne, 5~\% à la baisse, et a été suspendu de 09:25 à 09:30. Le plus grand
fonds indiciel a ouvert en baisse de 5.2~\%, était en baisse de 7.8~\% à 09:35, et a clôturé en baisse de 4.2~\%.
À 09:35, la plus grande bourse cotée n'avait ouvert que 38~\% de ses actions de l'indice
principal, si bien que pendant quelques minutes l'indice a été calculé en partie à partir
des cours de clôture du vendredi, et les arbitragistes des fonds n'avaient
pas de panier à valoriser (voir la \cref{prop:m1:exchange-traded-funds:stale}).
Il y a eu ce jour-là 1\,278 suspensions \omterm{def:m1:stress-case-studies:luld}{Limit Up--Limit Down (LULD)}~: 1\,058 dans 327 produits
cotés en bourse et 220 dans 144 autres titres. Un cinquième des produits cotés en
bourse a baissé de 20~\% ou plus à un moment donné, contre 5~\% des actions ordinaires,
bien que les quatre cinquièmes d'entre eux n'aient jamais été suspendus.

Le mécanisme conçu après 2010 a fonctionné comme prévu et a produit une nouvelle
défaillance. Un fonds dont les \omterm{def:m1:what-a-trading-firm-does:market-maker}{teneurs de marché} ne peuvent pas valoriser le panier n'a
pas d'ancrage~; une suspension arrête la négociation sans fournir l'ancrage~; et une
enchère de réouverture dans un fonds sans ancrage, avec des ordres stop en file
d'attente, peut le rouvrir à la bande puis le suspendre à nouveau. Les suspensions répétées dans des produits dont les actions sous-jacentes
se négociaient normalement ont été la signature de la journée.

\begin{omfigure}
\begin{tikzpicture}
\begin{axis}[width=9.5cm,height=4.4cm,ybar,bar width=16pt,
  ylabel={nombre},
  symbolic x coords={exchange-traded products,other securities},xtick={exchange-traded products,other securities},xticklabels={{produits cotés},{autres titres}},
  tick label style={font=\small},label style={font=\small},xticklabel style={font=\small\strut},
  ymin=0,ymax=1300,ymajorgrids,grid style={gray!25},enlarge x limits=0.5,
  scaled y ticks=false,yticklabel style={/pgf/number format/1000 sep={\,}},
  nodes near coords,nodes near coords style={font=\footnotesize,fill=white,inner sep=1pt,/pgf/number format/1000 sep={\,}},
  legend columns=2,legend style={font=\footnotesize,at={(0.5,-0.32)},anchor=north,draw=none,fill=none,/tikz/every even column/.append style={column sep=8pt}},
  table/col sep=comma]
\addplot[fill=omThm!60,draw=omThm] table[x=group,y=halts]{figdata/markets-1/31-stress-case-studies/halts.csv};
\addplot[fill=omDef!60,draw=omDef] table[x=group,y=securities]{figdata/markets-1/31-stress-case-studies/halts.csv};
\legend{suspensions,titres suspendus}
\end{axis}
\end{tikzpicture}

\omcaption{Suspensions Limit Up--Limit Down du 24 août 2015. Les fonds qui ont
été suspendus l'ont été en moyenne trois fois chacun. Données~: la note de
recherche du régulateur, décembre 2015.}\label{fig:m1:stress-case-studies:halts}
\end{omfigure}

\section{5 février 2018 et janvier 2021}

Deux épisodes ont déjà été évoqués. Le 5 février 2018
(\cref{ex:m1:volatility-first-contact:feb2018}), des produits promettant
l'inverse, ou un multiple, de la volatilité, devaient acheter ces contrats à terme
sur la volatilité à la clôture, en proportion de la hausse de la journée~; leurs achats étaient la hausse elle-même, et le
plus grand produit inverse a perdu presque toute sa valeur en une séance. Le flux
était mécanique, public et concentré en quelques minutes~: le rééquilibrage
suivait l'arithmétique de la \cref{prop:m1:exchange-traded-funds:rebalance}, appliquée à
un marché trop petit pour la supporter.

En janvier 2021 (\cref{ex:m1:stock-loan-and-short-selling:gme}), une action
dont l'\omterm{def:m1:stock-loan-and-short-selling:si}{encours vendeur à découvert} atteignait 123~\% de son flottant a progressé de 2\,700~\% en trois semaines. Les
mécanismes mis à rude épreuve n'étaient pas cotés sur une bourse. Les \omterm{def:m1:what-a-trading-firm-does:market-maker}{teneurs de marché} sur options ont vendu
les calls achetés par des particuliers, en se couvrant par l'achat d'actions
(\cref{def:m1:zero-day-options-and-retail-flow:dealer})~; les vendeurs à découvert se sont rachetés~;
et le modèle de marge de la chambre de compensation, voyant la volatilité et les
positions unilatérales des clients des courtiers de détail, a relevé leurs
exigences de plusieurs fois avant l'aube (\cref{ch:m1:clearing-and-settlement}, \cref{ch:m1:margin}). Certains courtiers
n'ont pas pu ou pas voulu financer l'\omterm{def:m1:financing:margincall}{appel de marge} et ont restreint leurs clients
aux seuls rachats de clôture. Le rapport du personnel a établi que la couverture ne représentait qu'une faible part des
achats et que la hausse était portée par le sentiment~; ce que l'épisode a ajouté
à la liste des facteurs susceptibles d'arrêter un marché, c'est le bilan d'un courtier
de détail.

\begin{definition}[Spirale de liquidité]\label{def:m1:stress-case-studies:spiral}
Une \emph{spirale de liquidité}\index{spirale de liquidité} est une boucle dans laquelle la baisse
des cours réduit la capacité des intermédiaires à détenir des positions (par
les pertes, les \omterm{def:m1:financing:margincall}{appels de marge}, les limites de risque ou l'incertitude sur la validité de leurs
transactions), leur retrait réduit la profondeur, et une profondeur réduite fait que le même
flux déplace davantage les cours.
\end{definition}

\begin{definition}[Transaction manifestement erronée]\label{def:m1:stress-case-studies:erroneous}
Une \emph{transaction manifestement erronée}\index{transaction manifestement erronée} est une transaction qu'une
bourse peut l'annuler a posteriori parce que son prix était trop éloigné
d'un cours de référence. Le 6 mai 2010, le seuil, 60~\%, a été fixé après
coup~; les bandes de prix existent pour que de telles transactions ne se produisent
pas en premier lieu.
\end{definition}

\section{Ce que les quatre cas ont en commun}

\begin{method}[Lire un incident]\label{met:m1:stress-case-studies:read}
\begin{enumerate}
\item \textbf{Le flux.} Qui devait négocier, combien, selon quelle règle et avant quand~?
Dans chaque cas, le flux était insensible au prix~: un algorithme suivant le volume,
des ordres stop-loss, un rééquilibrage quotidien, un \omterm{def:m1:financing:margincall}{appel de marge}.
\item \textbf{Les absorbeurs.} Qui prend normalement la contrepartie, et ce qui les limite~:
l'inventaire, le capital, une couverture qui doit exister, la certitude que
les transactions tiendront~?
\item \textbf{La boucle.} Par quelle variable le mouvement de prix revient-il
dans le flux ou dans la capacité des absorbeurs~: le volume, la volatilité, la marge, une
valorisation périmée~?
\item \textbf{Le circuit.} Qu'est-ce qui l'a arrêté~: une pause, une bande, une
enchère, un changement de règle, de nouveaux capitaux~? Qu'a cassé le mécanisme d'arrêt lui-même~?
\item \textbf{Les suites.} Quelles transactions ont tenu~? Qui a supporté la perte~? Quelle
règle a changé~?
\end{enumerate}
\end{method}

\begin{center}\small
\begin{tabular}{p{2.2cm}p{3.3cm}p{3.4cm}p{3.6cm}}
\toprule
 & Flux insensible au prix & Pourquoi les absorbeurs se sont retirés & Ce qui a changé ensuite \\
\midrule
6 mai 2010 & algorithme de vente suivant le volume & limites d'inventaire~; crainte d'annulation de transactions & \omterm{def:m1:stress-case-studies:flash}{cotations fictives} interdites (novembre 2010)~; bandes de prix \\
24 août 2015 & ordres au marché et ordres stop à l'ouverture & aucun cours de panier pour les fonds~; ouvertures tardives & une enchère de réouverture harmonisée, avec des bandes élargies (2017) \\
5 février 2018 & rééquilibrage quotidien des produits sur la volatilité & marché des contrats à terme trop petit pour le flux & le plus grand produit inverse racheté par son émetteur \\
janvier 2021 & achats de calls par les particuliers~; rachats de positions vendeuses & marge de compensation des courtiers & \omterm{def:m1:clearing-and-settlement:cycle}{cycle de règlement-livraison} raccourci (\cref{ch:m1:clearing-and-settlement}) \\
\bottomrule
\end{tabular}
\end{center}

Pour une société de trading, les leçons sont opérationnelles. Il faut savoir lesquelles de ses
stratégies sont des absorbeurs et ce qui fait cesser chacune d'entre elles~; décider à l'avance
si elle s'arrête en élargissant, en réduisant ou en se coupant, car en cas d'événement
il ne reste que quatre minutes. Il faut savoir quelles règles de bourse déterminent si
ses transactions sont maintenues. Conserver la trésorerie nécessaire à l'\omterm{def:m1:financing:margincall}{appel de marge} qui accompagne
l'opportunité (\cref{pb:m1:margin:1}). Et tenir un enregistrement rejouable
(\cref{ch:m1:reading-market-data})~: une journée qui peut être rejouée est une journée
dont on peut tirer des enseignements.

\begin{omfigure}
\begin{tikzpicture}
\begin{axis}[width=10cm,height=4.6cm,
  xlabel={plus forte baisse tolérée (\%)},ylabel={profondeur nécessaire (milliers)},
  tick label style={font=\small},label style={font=\small},
  xmin=1.5,xmax=8.5,ymin=0,ymax=190,grid=major,grid style={gray!25},
  nodes near coords,nodes near coords style={font=\footnotesize,above=2pt,fill=white,inner sep=1pt,/pgf/number format/.cd,fixed,precision=0},
  table/col sep=comma]
\addplot[omDef,very thick,mark=*] table[x=max_fall_pct,y=depth_k]{figdata/markets-1/31-stress-case-studies/depth_needed.csv};
\draw[thin,dashed] (axis cs:1.5,100) -- (axis cs:8.5,100);
\node[font=\footnotesize,anchor=north east] at (axis cs:8.4,98) {profondeur réelle};
\end{axis}
\end{tikzpicture}

\omcaption{Dans le modèle jouet, la profondeur initiale qui aurait permis de contenir la
baisse du programme dans une taille donnée. Avec les rétroactions, 16~\% de profondeur
supplémentaire transforme une baisse de 9.8~\% en une baisse de 5~\%, et 61~\% de plus la ramène
à 2~\%, proche de la baisse sans aucune rétroaction~: la stabilité est très non linéaire en
profondeur, ce qui explique qu'elle disparaisse si brutalement. Données~: la
simulation du tutoriel.}\label{fig:m1:stress-case-studies:depth}
\end{omfigure}

\section{Tutoriel~: un retrait de liquidité}\label{tut:m1:stress-case-studies:withdrawal}

\begin{tutorial}
\textbf{Objectif.} Reproduire la forme d'un programme de vente suivant le volume dans un
carnet dont la profondeur dépend de la baisse déjà subie~; ajouter une pause~; trouver la
profondeur qui aurait absorbé le flux. \textbf{État final~:} les deux figures simulées de
ce chapitre.
\begin{tutsteps}
\item \textbf{Une minute du modèle,} en soixante étapes, afin que le résultat ne dépende
pas du pas de temps. Ventes à partir du volume de la minute précédente~; profondeur et
désistement à partir de la baisse par rapport à un cours moyen sur cinq minutes~; impact à
partir des ventes rapportées à la profondeur.
\omcode{code/markets-1/31-stress-case-studies/python/withdrawal.py}{20}{36}{Un programme de vente, un carnet qui se retire et une pause.}\label{lst:m1:stress-case-studies:run}
\item \textbf{Désactiver les rétroactions} une à une, puis ralentir le programme à 3~\% du
volume, et observer ce qui fait le plus de dégâts.
\item \textbf{Bandes et pauses} sous forme de machine à états sur un flux de transactions.
\omcode{code/firm/replay/firm_replay.py}{47}{77}{Cours de référence, état de limitation et pause.}\label{lst:m1:stress-case-studies:luld}
\end{tutsteps}
\textbf{Pour aller plus loin.} Donner au programme une \omterm{def:m1:asian-and-emerging-equity-markets:limit}{limite de variation de cours} (ne pas
vendre plus de 1~\% sous le cours d'arrivée) et comparer son coût et sa durée.
Puis ajouter une reprise~: des acheteurs qui reviennent en proportion de la distance sous une
fair value qui se déplace lentement.
\end{tutorial}

\section{Construction~: le harnais de rejeu d'incidents}\label{bld:m1:stress-case-studies:replay}

\begin{build}
\bfield{Rôle} La mini-société de trading fait passer chaque stratégie par une bibliothèque de
mauvaises journées avant de négocier, et rejoue ensuite ses propres incidents. Le harnais
transforme un flux de transactions en la chronologie dont une analyse post-mortem a besoin.
\bfield{Interface} \texttt{Tick(t, bid, ask, trade)}~; \texttt{Luld(band\_pct,
\ldots)} avec \texttt{pct(t)}, \texttt{bands(t)}, \texttt{on\_tick(tick)} et une
liste d'\texttt{events}~; \texttt{replay(ticks, luld)} qui renvoie l'état et les
bandes à chaque tick~; \texttt{depth\_within(levels, mid, bp)}.
\bfield{Règles} Les règles de bandes de l'\cref{dat:m1:stress-case-studies:luld}, paramétrées~:
fenêtre, mouvement minimal, durée minimale de maintien, périodes de doublement, durée de
limitation, durée de la pause. Après une pause, la première transaction fixe une nouvelle
référence. Les entrées proviennent du normaliseur de flux de
la \cref{bld:m1:reading-market-data:feed}, de sorte qu'une journée enregistrée puisse être
rejouée sans modification.
\bfield{Tests d'acceptation} \texttt{code/firm/replay/tests/}~: doublement à l'ouverture et à
la clôture~; une référence qui ne bouge que de 1~\% et après 30 secondes~; un état de
limitation qui devient une pause après 15 secondes, puis une réouverture~; un état de
limitation qui se résorbe à temps~; une profondeur proche du mid.
\bfield{Pour aller plus loin} Une bibliothèque de scénarios~: pour chacune des quatre
journées de ce chapitre, un flux de transactions synthétique reproduisant les
caractéristiques documentées, et un rapport sur ce que chacune des stratégies de la société aurait fait.
\end{build}

\begin{omsources}
\item US Commodity Futures Trading Commission et US Securities and Exchange
Commission, \emph{Findings Regarding the Market Events of May 6, 2010}, 30
septembre 2010.
\item US Securities and Exchange Commission, Division of Trading and Markets,
\emph{Research Note: Equity Market Volatility on August 24, 2015}, décembre
2015.
\item Bank for International Settlements, «~The equity market turbulence of 5
February: the role of exchange-traded volatility products~», \emph{BIS
Quarterly Review}, mars 2018.
\item US Securities and Exchange Commission, \emph{Staff Report on Equity and
Options Market Structure Conditions in Early 2021}, 14 octobre 2021.
\item M.~Brunnermeier et L.~H.~Pedersen, «~Market liquidity and funding
liquidity~», \emph{Review of Financial Studies} 22 (2009).
\end{omsources}

\section{Exercices}

\begin{exercise}[$\star$]\label{exo:m1:stress-case-studies:1}
Le programme E-mini du 6 mai 2010 portait sur 75\,000 contrats, soit 4.1~milliards de dollars.
Donner le niveau implicite de l'indice (multiplicateur 50~\$). Entre 14 h 32 et 14 h 45, il a vendu
environ 35\,000 contrats~: quelle part du programme cela représente-t-il,
et à quel rythme moyen par minute~?
\end{exercise}

\begin{exercise}[$\star$]\label{exo:m1:stress-case-studies:2}
Le cours de référence d'une action est de 40.00 et sa bande de 5~\%. Donner les bandes à 11:00
et à 09:40. Le meilleur cours vendeur tombe à 38.00 à 11:00:00 et s'y maintient.
Quand la pause est-elle déclarée, et jusqu'à quand dure-t-elle au moins~?
\end{exercise}

\begin{exercise}[$\star$]\label{exo:m1:stress-case-studies:3}
D'après la \cref{fig:m1:stress-case-studies:halts}~: combien de suspensions par titre suspendu,
pour les produits cotés en bourse et pour les autres titres~?
Qu'est-ce que l'écart suggère~?
\end{exercise}

\begin{exercise}[$\star\star$]\label{exo:m1:stress-case-studies:4}
Un algorithme vend 9~\% du volume de la minute précédente. Le volume est de 20\,000
contrats par minute en conditions calmes. Combien de temps faut-il pour vendre 75\,000
contrats~? Si la vente propre de l'algorithme et le désistement qu'elle provoque triplent
le volume, combien de temps cela prend-il~? Pourquoi «~un pourcentage du volume~» n'est-il pas une
consigne neutre~?
\end{exercise}

\begin{exercise}[$\star\star$]\label{exo:m1:stress-case-studies:5}
Un \omterm{def:m1:what-a-trading-firm-does:market-maker}{teneur de marché} est tenu de coter à double sens en permanence. Expliquer comment une
\omterm{def:m1:stress-case-studies:flash}{cotation fictive} respecte la lettre de cette obligation, ce qu'il est advenu des ordres au
marché qui ont rencontré des \omterm{def:m1:stress-case-studies:flash}{cotations fictives} le 6 mai 2010, et ce qui a remplacé les \omterm{def:m1:stress-case-studies:flash}{cotations fictives}.
\end{exercise}

\begin{exercise}[$\star\star$]\label{exo:m1:stress-case-studies:6}
Le 24 août 2015, un fonds indiciel a été négocié 20~\% en dessous de la valeur de ses
avoirs, alors que les marchés de ces avoirs fonctionnaient. En utilisant la
\cref{prop:m1:exchange-traded-funds:band}, dire quel terme de la bande d'arbitrage
était devenu très grand, et pourquoi une pause de cinq minutes ne l'a pas réduit.
\end{exercise}

\begin{exercise}[$\star\star\star$]\label{exo:m1:stress-case-studies:7}
\emph{Codage.} Avec \texttt{run} et \texttt{trough}, reproduire les trois
chutes de la \cref{fig:m1:stress-case-studies:paths}. Puis désactiver seulement
le désistement, puis seulement le retrait, et exécuter le modèle complet avec
une participation de 3~\%. Pourquoi chaque rétroaction a-t-elle besoin de l'autre pour être dangereuse~?
\end{exercise}

\begin{exercise}[$\star\star\star$]\label{exo:m1:stress-case-studies:8}
\emph{Trouver la faille.} Un rapport de risque~: «~Notre stratégie de \omterm{def:m1:what-a-trading-firm-does:market-maker}{teneur de
marché} est sûre en cas de krach~: elle est à plat à la fin de chaque journée,
son pire jour sur cinq ans de backtest représente $-0.4\%$ du capital, et elle
élargit automatiquement ses cotations quand la volatilité augmente.~» Nommer trois
choses que les quatre jours du chapitre ont faites et que ce rapport ne peut pas voir.
\end{exercise}

\section{Problème~: Cinq secondes}

\begin{problem}[{Problème du week-end --- quelle profondeur aurait suffi~?}]\label{pb:m1:stress-case-studies:1}
Utiliser le modèle jouet du chapitre avec ses paramètres par défaut~: profondeur
initiale côté acheteur de 100\,000 contrats, volume calme de 20\,000 par minute,
un programme de 75\,000 contrats vendant 9~\% du volume de la minute précédente,
impact de 3~\% pour la vente d'une quantité égale à la profondeur courante. Avec
$f$ la baisse du cours en dessous de sa moyenne sur les cinq dernières minutes
environ, la profondeur est multipliée par $0.1 + 0.9\exp(-150 f)$ et le volume calme par $1 + 100 f$.

\medskip\noindent\textbf{Partie I --- La première minute.}
\begin{enumerate}
\item Donner les ventes du programme pendant la première minute.
\item Donner la variation de cours qu'elles provoquent si la profondeur reste à son niveau initial.
\item Lorsque le cours est inférieur de 1~\% à sa moyenne récente, donner la profondeur et le
volume.
\item De combien la profondeur peut-elle devenir petite, et à quelle baisse récente se trouve-t-elle
à un dixième près de ce plancher~?
\item Pour une baisse récente de 1~\%, combien le programme vend-il en une minute,
et que coûte cette minute~?
\end{enumerate}

\medskip\noindent\textbf{Partie II --- Le chemin complet.}
\begin{enumerate}[resume]
\item La simulation donne des baisses de 2.2~\% avec un carnet figé et de 9.8~\% avec les
rétroactions. Combien de fois la seconde est-elle plus coûteuse~?
\item Évaluer l'écart sur 75\,000 contrats à 54\,667~\$ chacun, en prenant
la vente moyenne à la moitié du point bas.
\item Avec une pause à $-5\%$ après laquelle les acheteurs mesurent la baisse à partir du
cours de réouverture, le point bas est à 7.0~\%. Qu'a économisé la pause~?
\item Dans le modèle, pourquoi la pause fonctionne-t-elle~? La rattacher à ce que le rapport officiel
dit qu'il s'est passé pendant les cinq secondes.
\end{enumerate}

\medskip\noindent\textbf{Partie III --- La profondeur qui manquait.}
\begin{enumerate}[resume]
\item D'après la \cref{fig:m1:stress-case-studies:depth}, donner la profondeur initiale qui maintient la baisse
sous 2~\%, puis sous 5~\%.
\item Exprimer chacune comme un multiple des 100\,000 réels.
\item Pourquoi la relation est-elle si non linéaire~?
\item Le programme aurait pu se voir imposer un taux de participation
de 3~\%. Sans exécuter le modèle, que changerait-il, et que coûterait-il
au vendeur~?
\item Quelle partie était la mieux placée pour empêcher l'événement au moindre coût~?
\end{enumerate}

\medskip\noindent\textbf{Partie IV --- Jugement.}
\begin{enumerate}[resume]
\item Les intermédiaires ont échangé 27\,000 contrats en quatorze secondes et en ont acheté 200 en net.
Fournissaient-ils de la liquidité~?
\item Les 20\,000 transactions à des prix absurdes auraient-elles dû être annulées~? Donner
l'argument de chacun des deux côtés.
\item Les bandes et les pauses existent désormais. Citer une chose qu'elles ont corrigée et une qu'elles n'ont pas corrigée,
à partir du 24 août 2015.
\item Que doit toujours posséder un algorithme d'exécution qui suit le volume~?
\item Énoncer le \emph{résultat nommé}~: la profondeur initiale qui aurait absorbé
le programme avec une baisse d'au plus 2~\%.
\item En une phrase~: qu'est-ce que la liquidité, le jour où elle compte~?
\end{enumerate}
\end{problem}

\section{Questions d'entretien}

\begin{interviewq}[$\star$ \iqroles{trader, chercheur, développeur}]\label{iq:m1:stress-case-studies:1}
Que s'est-il passé le 6 mai 2010~?
\end{interviewq}

\begin{interviewq}[$\star$ \iqroles{trader, développeur}]\label{iq:m1:stress-case-studies:2}
Comment fonctionne le Limit Up--Limit Down~?
\end{interviewq}

\begin{interviewq}[$\star\star$ \iqroles{trader, chercheur}]\label{iq:m1:stress-case-studies:3}
Pourquoi les \omterm{def:m1:exchange-traded-funds:etf}{fonds indiciels cotés} (ETF) se sont-ils échangés si loin en dessous de la valeur de leurs
avoirs le 24 août 2015~?
\end{interviewq}

\begin{interviewq}[$\star\star$ \iqroles{chercheur, ingénieur ML}]\label{iq:m1:stress-case-studies:4}
Votre backtest couvre cinq années calmes. Comment évaluez-vous le
comportement de la stratégie dans un krach~?
\end{interviewq}

\begin{interviewq}[$\star\star$ \iqroles{trader, développeur}]\label{iq:m1:stress-case-studies:5}
Votre système de \omterm{def:m1:what-a-trading-firm-does:market-maker}{teneur de marché} voit le cours chuter de 3~\% en deux minutes sur un volume en hausse.
Que doit-il faire, et qui décide~?
\end{interviewq}

\begin{interviewq}[$\star\star\star$ \iqroles{chercheur, trader, banque}]\label{iq:m1:stress-case-studies:6}
Qu'ont en commun les krachs provoqués par la structure de marché, et que
chercheriez-vous pour anticiper le prochain~?
\end{interviewq}
